Binary operators get one surrounding space, unary signs and scripts stay tight:
$a + 2 * 1^5$, $f(x) = x + 1$, $a/b + c$, $a = b = c$, and $a <= b$.

Slash adjacency is symmetric: $a/b$, $c / d$, $e / f$, and $g / h$.

Formatter-introduced operator gaps also make slash adjacency symmetric:
$a / \sim$, $b \sim / c$, and $d / = e$.

Unary is glued: $-x + 1$, $x = -b$, $2 * -1$, and $2^{-5}$.

A sign after an opening delimiter is unary: $(-1)^n$, $[-x]$, $f(-1)$, and
$\{-y\}$, but stays binary after an operand: $(x - 1)$.

Command operators are spaced too: $a \cdot b$, $a \times b$, and $x \leq y$.

Generated Unicode and command classes use the same policy: $a ≤ b$, $a ⊕ b$, and
$x \nleq y$.

Colon-prefixed relations stay indivisible: $a := b$, $a := b$, $a ::= b$, and
$a : = b$; scripted forms stay indivisible too: $a :=_i b$ and $a ::=^j c$; an
ordinary colon remains punctuation: $f:x$.

Trivia keeps a scripted equals separate from its colon: $a : =_k b$.

Script punctuation stays tight while control-word operators stay readable: $x^{a+b}$,
$\sum_{i=1}^{n / 2} a_i$, $x_{a \in A}$, and
$x_{\frac{a+b}{c \leq d}}$.

A number in scientific notation is not special-cased: $1e - 5$.
